\section*{Phase-2 Tier-3 Plan: Robustness and Generalization Study}

\section*{Supernova Classification Project}

\section{Objective}

Phase-2 Tier-1 demonstrated that a compact set of interpretable photometric features
can preserve nearly all classification performance for Type Ia supernova identification
while improving interpretability and computational efficiency.

Phase-2 Tier-2 analyzed the compact representation through controlled ablation
experiments and identified which features and feature groups contribute most strongly
to classification performance.

The objective of Phase-2 Tier-3 is to test whether the compact feature representation
captures real astrophysical information rather than dataset-specific correlations.
This phase evaluates the robustness, stability, and generalization of the compact
feature set across different models, data splits, and observational conditions.

The central question addressed in this phase is:

\textit{Does the compact feature representation remain reliable under changes in model,
data sampling, and noise conditions, indicating that it encodes genuine physical
information from supernova light curves?}

\section{Scientific Questions}

Phase-2 Tier-3 focuses on the following questions:

\begin{enumerate}
\item Do the same compact features remain important across different classifiers?
\item Is the classification performance stable under different train-test splits?
\item How sensitive is the compact representation to noise or missing observations?
\item Do different feature-importance methods identify the same dominant features?
\item Can the minimal core feature subset discovered in Tier-2 generalize to new conditions?
\item Do the retained features correspond to known physical properties of supernova light curves?
\end{enumerate}

\section{Tier-1 Reference Baseline}

All Tier-3 experiments are evaluated relative to the frozen compact Tier-1 baseline.

\begin{center}
\begin{tabular}{l c}
\textbf{Metric} & \textbf{Value} \\
Number of features & 16 \\
F1 score & 0.8442 \\
ROC-AUC & 0.9766 \\
PR-AUC & 0.9278 \\
\end{tabular}
\end{center}

\section{Compact Feature Representation}

The compact feature set obtained in Tier-1 consists of sixteen interpretable
photometric features organized into four physical groups.

\subsection{Brightness Features}
\begin{itemize}
\item r mean flux
\item g mean flux
\item z peak flux
\item i peak flux
\item r peak flux
\end{itemize}

\subsection{Color Features}
\begin{itemize}
\item peak color g minus r
\item peak color r minus i
\item peak color i minus z
\end{itemize}

\subsection{Variability Features}
\begin{itemize}
\item i std flux
\item z std flux
\item r std flux
\item i amplitude
\end{itemize}

\subsection{Temporal Features}
\begin{itemize}
\item r time of peak
\item i time of peak
\item z time of peak
\item time span
\end{itemize}

\section{Experiments}

\subsection{Experiment A: Model Robustness}

The compact feature set is evaluated using multiple classifiers:
\begin{itemize}
\item XGBoost (baseline)
\item Random Forest
\item Logistic Regression
\item Support Vector Machine
\end{itemize}

For each model the following metrics are recorded:
\begin{itemize}
\item F1 score
\item ROC-AUC
\item PR-AUC
\end{itemize}

\subsection{Experiment B: Cross-Validation Stability}

Instead of a single train-test split, repeated cross-validation is performed:
\begin{itemize}
\item k-fold cross validation
\item repeated random splits
\end{itemize}

For each run the mean and standard deviation of the metrics are recorded.

\subsection{Experiment C: Noise and Missing Data Tests}

Artificial perturbations are introduced to simulate realistic survey conditions:
\begin{itemize}
\item reduced number of observations
\item added noise to flux values
\item removal of selected bands
\item shortened time coverage
\end{itemize}

Performance is compared with the baseline.

\subsection{Experiment D: Feature Importance Consistency}

Feature importance is computed using different methods:
\begin{itemize}
\item model gain importance
\item permutation importance
\item SHAP values
\item ablation results from Tier-2
\end{itemize}

\subsection{Experiment E: Minimal Core Generalization}

The minimal feature subsets identified in Tier-2 are tested again under
new conditions using candidate subset sizes:
\begin{itemize}
\item top 5 features
\item top 8 features
\item top 10 features
\end{itemize}

\subsection{Experiment F: Physical Interpretation Check}

The retained features are interpreted in terms of known supernova physics:
\begin{itemize}
\item color features $\rightarrow$ temperature evolution
\item peak flux $\rightarrow$ luminosity scale
\item time of peak $\rightarrow$ light-curve width
\item variability $\rightarrow$ shape of the light curve
\end{itemize}

\section{Evaluation Metrics}

The following metrics are used:
\begin{itemize}
\item F1 score
\item Precision-Recall AUC
\item ROC-AUC
\end{itemize}

F1 and PR-AUC are emphasized because they better reflect performance under
class imbalance.

\section{Repository Structure}

New Tier-3 files:
\begin{verbatim}
phase2_tier3_model_compare.py
phase2_tier3_cv_stability.py
phase2_tier3_noise_test.py
phase2_tier3_importance_compare.py
phase2_tier3_minimal_generalization.py
phase2_tier3_summary.py
phase2_tier3_plan.tex
\end{verbatim}

Outputs:
\begin{verbatim}
results/phase2_tier3/
plots/phase2_tier3/
\end{verbatim}

\section{Workflow}

\begin{enumerate}
\item Freeze Tier-1 compact baseline
\item Run model robustness tests
\item Run cross-validation stability tests
\item Run noise and missing-data experiments
\item Compare feature-importance methods
\item Test minimal feature subsets
\item Interpret results physically
\item Document Tier-3 results
\end{enumerate}

\section{Final Goal}

Phase-2 Tier-3 determines whether the compact interpretable feature set remains
robust, stable, and physically meaningful under changes in model choice, data
sampling, and observational quality. This phase establishes the scientific
validity of the compact representation as more than a dataset-specific shortcut.
